
\section{$A_\infty$ 代数的同伦等价：代数框架}
\label{sec:4-2}

\subsection{$[0,1]\times C$ 的模型}
\label{subsec:4-2-1}

下一小节将定义两个（滤过）$A_\infty$ 同态彼此同伦。这个概念有多种
定义方式，也有多篇文献讨论它（\textbf{[Kel01, MSS02, Smi00]}）。为统一
这些不同途径，我们作如下处理。在代数拓扑中，两个映射
$f,f'\colon X\to Y$ 之间的同伦由映射
$H\colon[0,1]\times X\to Y$ 定义。因此，在（滤过）$A_\infty$ 代数的
语境中，我们先定义运算 $[0,1]\times-$——即与 $[0,1]$ 取乘积——的代数
类比。

令 $R$ 为含幺交换环，令 $\overline C$ 为 $R$ 上的 $A_\infty$ 代数，令
$C$ 为 $\Lambda_{0,\mathrm{nov}}(R)$ 上的滤过 $A_\infty$ 代数。这里
$\Lambda_{0,\mathrm{nov}}(R)$ 是 $R$ 上的通用诺维科夫环，见第 1 章
(Conv.4)。我们始终假定 $\overline C$ 与 $C$ 分别是 $R$ 与
$\Lambda_{0,\mathrm{nov}}(R)$ 上自由且可数生成的模，并且
$\overline C$ 是 $C$ 的 $R$-约化（定义 3.2.20）。以后在不会混淆时，
以 $\Lambda_{0,\mathrm{nov}}$ 代替 $\Lambda_{0,\mathrm{nov}}(R)$。

\begin{definition}{4.2.1}\label{def:4-2-1}
设 $\overline{\mathfrak C}$ 是一个 $A_\infty$ 代数，并带有
$A_\infty$ 同态
\[
 \overline{\mathrm{Incl}}\colon\overline C\longrightarrow
 \overline{\mathfrak C},\qquad
 \overline{\mathrm{Eval}}_{s=0}\colon\overline{\mathfrak C}
 \longrightarrow\overline C,\qquad
 \overline{\mathrm{Eval}}_{s=1}\colon\overline{\mathfrak C}
 \longrightarrow\overline C.
\]
若下列条件成立，就称 $\overline{\mathfrak C}$ 是
$[0,1]\times\overline C$ 的一个\emph{模型}\liftIndex{I141}\liftIndex{I143}：
\begin{enumerate}[label=\textup{(4.2.2.\arabic*)},leftmargin=5.4em]
 \item $\overline{\mathrm{Incl}}_k\colon B_k\overline C[1]
  \to\overline{\mathfrak C}[1]$ 在 $k\ne1$ 时为零；
  $\overline{\mathrm{Eval}}_{s=0}$ 与
  $\overline{\mathrm{Eval}}_{s=1}$ 也一样。
 \item $\overline{\mathrm{Eval}}_{s=0}\circ
  \overline{\mathrm{Incl}}=
  \overline{\mathrm{Eval}}_{s=1}\circ
  \overline{\mathrm{Incl}}=\mathrm{id}$。
 \item $\overline{\mathrm{Incl}}_1\colon
  (\overline C,\overline{\mathfrak m})\to
  (\overline{\mathfrak C},\overline{\mathfrak m})$ 是上链同伦等价，且
  $(\overline{\mathrm{Eval}}_{s=0})_1$、
  $(\overline{\mathrm{Eval}}_{s=1})_1\colon
  (\overline{\mathfrak C},\overline{\mathfrak m})\to
  (\overline C,\overline{\mathfrak m})$ 也是上链同伦等价。
 \item 上链同态
  $(\overline{\mathrm{Eval}}_{s=0})_1\oplus
  (\overline{\mathrm{Eval}}_{s=1})_1\colon
  \overline{\mathfrak C}\to\overline C\oplus\overline C$ 是满射。
\end{enumerate}
若 $\overline C$、$\overline{\mathfrak C}$、
$\overline{\mathrm{Incl}}$、$\overline{\mathrm{Eval}}_{s=0}$ 与
$\overline{\mathrm{Eval}}_{s=1}$ 都含幺，就称
$\overline{\mathfrak C}$ 是 $[0,1]\times\overline C$ 的
\emph{含幺模型}。

设 $\mathfrak C$ 是一个滤过 $A_\infty$ 代数，并带有滤过
$A_\infty$ 同态
\[
 \mathrm{Incl}\colon C\longrightarrow\mathfrak C,
 \qquad
 \mathrm{Eval}_{s=0}\colon\mathfrak C\longrightarrow C,
 \qquad
 \mathrm{Eval}_{s=1}\colon\mathfrak C\longrightarrow C.
\]
若下列条件成立，就称 $\mathfrak C$ 是 $[0,1]\times C$ 的一个模型：
\begin{enumerate}[label=\textup{(4.2.3.\arabic*)},leftmargin=5.4em]
 \item $\mathrm{Incl}_k\colon B_kC[1]\to\mathfrak C[1]$ 在
  $k\ne1$ 时为零；$\mathrm{Eval}_{s=0}$ 与
  $\mathrm{Eval}_{s=1}$ 也一样。\translatornote{原书此处把
  $\mathrm{Incl}_k$ 的靶写成 $\mathfrak C$，漏去移位 $[1]$；同一定义的
  无滤过版本以及 $A_\infty$ 同态分量的型别均要求靶为
  $\mathfrak C[1]$，译文据此补正。}
 \item $\mathrm{Eval}_{s=0}\circ\mathrm{Incl}=
  \mathrm{Eval}_{s=1}\circ\mathrm{Incl}=\mathrm{id}$。
 \item $\mathrm{Incl}_1$ 在复形之间诱导上链同伦等价
  $(\overline C,\overline{\mathfrak m})\to
  (\overline{\mathfrak C},\overline{\mathfrak m})$；
  $(\mathrm{Eval}_{s=0})_1$ 与 $(\mathrm{Eval}_{s=1})_1$ 在复形之间
  诱导上链同伦等价
  $(\overline{\mathfrak C},\overline{\mathfrak m})\to
  (\overline C,\overline{\mathfrak m})$。
 \item 同态
  $(\mathrm{Eval}_{s=0})_1\oplus(\mathrm{Eval}_{s=1})_1\colon
  \mathfrak C\to C\oplus C$ 是满射。
\end{enumerate}
若 $C$、$\mathfrak C$、$\mathrm{Incl}$、$\mathrm{Eval}_{s=0}$ 与
$\mathrm{Eval}_{s=1}$ 都含幺，就称 $\mathfrak C$ 是
$[0,1]\times C$ 的含幺模型。
\end{definition}
\liftIndex{I142}\liftIndex{I262}

\begin{remark}{4.2.4}\label{rem:4-2-4}
\begin{enumerate}[label=\textup{(\arabic*)},leftmargin=2.8em]
 \item 在大多数情形下（例如 $R$ 是域或 $\mathbb Z$），条件
  (4.2.2.3) 等价于要求 $\overline{\mathrm{Incl}}_1$、
  $(\overline{\mathrm{Eval}}_{s=0})_1$ 与
  $(\overline{\mathrm{Eval}}_{s=1})_1$ 在
  $\overline{\mathfrak m}_1$ 上同调上诱导同构；条件 (4.2.3.3) 也有
  类似说明。注意我们假定 $\overline C$ 是自由可数生成的 $R$-模，故
  (4.2.2.4) 中的映射
  $(\overline{\mathrm{Eval}}_{s=0})_1\oplus
  (\overline{\mathrm{Eval}}_{s=1})_1$ 实际上是分裂满射。
 \item 条件 (4.2.3.3) 只要求把系数环
  $\Lambda_{0,\mathrm{nov}}$ 约化到 $R$ 以后得到上链同伦等价。这是因为
  $\mathfrak m_1\circ\mathfrak m_1$ 可能不为零。
 \item 对下面将构造的模型，有
  $\mathrm{Incl}_1=\overline{\mathrm{Incl}}_1$、
  $\mathrm{Eval}_1=\overline{\mathrm{Eval}}_1$；我们不把这一点列入定义。
\end{enumerate}
\end{remark}

\begin{example}{4.2.5}\label{ex:4-2-5}
令 $X$ 为光滑流形。取 $\overline C=\Lambda(X)$ 为 $X$ 上光滑微分形式的
德拉姆复形，取 $\overline{\mathfrak C}=\Lambda(\mathbb R\times X)$ 为
$\mathbb R\times X$ 上光滑微分形式的德拉姆复形。令
$I_i\colon X\to\mathbb R\times X$（$i=0,1$）为嵌入
$I_i(x)=(i,x)$，令 $\pi\colon\mathbb R\times X\to X$ 为投影。取
$\overline{\mathrm{Incl}}_1=\pi^*$、
$(\overline{\mathrm{Eval}}_{s=i})_1=I_i^*$，则
$\overline{\mathfrak C}$ 显然满足条件 (4.2.2.1)--(4.2.2.4)。
\end{example}

本小节余下部分将对 $\Lambda_{0,\mathrm{nov}}(R)$ 上任意滤过
$A_\infty$ 代数 $C$ 构造 $[0,1]\times C$ 的模型。先在
$R\supseteq\mathbb Q$ 时给出一个例子；这个模型适用于我们的绝大多数
目的。然后再对一般的 $R$ 构造 $[0,1]\times C$ 的模型，由此便能在一般
$R$（例如 $\mathbb Z$ 或 $\mathbb Z/2\mathbb Z$）上定义同伦等价。

先讨论 $R\supseteq\mathbb Q$ 的情形。下面的构造是
\textbf{[GrMo81]} 第 X 章中微分分次代数模型的 $A_\infty$ 类比。令
$(C,\mathfrak m)$ 为滤过 $A_\infty$ 代数，考虑“形式幂级数”
\begin{equation}
 x(s)=\sum_i x_i(s)T^{\lambda_i},
 \qquad
 y(s)=\sum_i y_i(s)T^{\lambda_i},
\tag{4.2.6}\label{eq:4-2-6}
\end{equation}
其中 $\lambda_i\geq0$ 且 $\lim_{i\to\infty}\lambda_i=\infty$。注意
$\lambda_i$ 是实数，不必是整数；这正是上面给“形式幂级数”加引号的原因。
假定系数函数
\[
 x_i\colon[0,1]\longrightarrow C^k[1]=C^{k+1},
 \qquad
 y_i\colon[0,1]\longrightarrow C^{k-1}[1]=C^k
\]
满足下列条件之一：
\begin{enumerate}[label=\textup{(4.2.7.\arabic*)},leftmargin=5.4em]
 \item $x_i(s)$、$y_i(s)$ 都是变量 $s$ 的多项式；
 \item $x_i(s)$、$y_i(s)$ 都是 $s\in[0,1]$ 的光滑函数；
 \item $x_i(s)$ 是 $s\in[0,1]$ 的连续分段多项式函数，$y_i(s)$ 是
  $s\in[0,1]$ 的分段多项式函数，但不要求连续；
 \item $x_i(s)$ 是 $s\in[0,1]$ 的连续分段光滑函数，$y_i(s)$ 是
  $s\in[0,1]$ 的分段光滑函数，但不要求连续。
\end{enumerate}
在 (4.2.7.2)、(4.2.7.4) 两种情形中，假定 $R=\mathbb R$ 或
$\mathbb C$。

\begin{remark}{4.2.8}\label{rem:4-2-8}
除能量滤过所诱导的拓扑外，我们不在向量空间 $C$ 上另取任何特定拓扑。
上面所谓光滑，是指相对于一组基的每个系数都光滑。
\end{remark}

\begin{definition}{4.2.9}\label{def:4-2-9}
定义分次模 $\operatorname{Poly}([0,1],C)$\liftIndex{I362}，使
\[
 \operatorname{Poly}([0,1],C)^{k+1}
 =\operatorname{Poly}([0,1],C)[1]^k
 \cong\operatorname{Poly}([0,1],C[1])^k
\]
是满足 (4.2.7.1) 的所有有序对 $(x,y)$ 所成的集合。分别把条件
(4.2.7.1) 换成 (4.2.7.2)、(4.2.7.3)、(4.2.7.4)，类似定义模
$C^\infty([0,1],C)$、$\operatorname{PPoly}([0,1],C)$ 与
$PC^\infty([0,1],C)$。

再令 $\overline C$ 为 $A_\infty$ 代数。定义分次模
$\overline{\operatorname{Poly}}([0,1],\overline C)$，使
\[
 \overline{\operatorname{Poly}}([0,1],\overline C)^{k+1}
 =\overline{\operatorname{Poly}}([0,1],\overline C)[1]^k
 \cong\overline{\operatorname{Poly}}([0,1],\overline C[1])^k
\]
是所有有序对 $(x,y)$ 所成的集合，其中 $x=x(s)$ 是取值于
$\overline C[1]^k=\overline C^{k+1}$ 的多项式，$y=y(s)$ 是取值于
$\overline C[1]^{k-1}=\overline C^k$ 的多项式。类似定义
$\overline{C^\infty}([0,1],\overline C)$、
$\overline{\operatorname{PPoly}}([0,1],\overline C)$ 与
$\overline{PC^\infty}([0,1],\overline C)$。
\end{definition}

\begin{remark}{4.2.10}\label{rem:4-2-10}
在 $\operatorname{Poly}([0,1],C)^k$ 的定义中，并不要求 $x_i(s)$ 作为
$s$ 的多项式，其次数关于 $i$ 一致有界。因此，若把
\eqref{eq:4-2-6} 改写为
\begin{equation}
 x(s)=\sum_{k=0}^{\infty}x_{(k)}(T)s^k,
\tag{4.2.11}\label{eq:4-2-11}
\end{equation}
其中 $x_{(k)}(T)\in C$，则右端不是有限和。不过，存在满足
$\lim_{k\to\infty}\lambda_k=\infty$ 的 $\lambda_k$，使
$x_{(k)}(T)\in F^{\lambda_k}C$。换言之，$x(s)$ 属于严格收敛幂级数环\liftIndex{I371}
$\Lambda_{0,\mathrm{nov}}\langle s\rangle$（见 \textbf{[BGR84]}
1.4.1 定义 1）。
\end{remark}

下面在 $\operatorname{Map}([0,1],C)$ 上定义滤过 $A_\infty$ 代数结构。
从现在起，$\operatorname{Map}$ 代表
$\operatorname{Poly}$、$C^\infty$、$\operatorname{PPoly}$ 或
$PC^\infty$ 之一：
\[
 \operatorname{Map}=\operatorname{Poly},\ C^\infty,
 \ \operatorname{PPoly}\ \text{或}\ PC^\infty.
\]
先定义 $\mathfrak m_k$（$k\geq1$）。令
$\mathbf x^i=(x^{(i)},y^{(i)})\in
\operatorname{Map}([0,1],C)^{d_i}$，其中 $(x^{(i)},y^{(i)})$ 如定义
4.2.9；特别地，$d_i$ 是 $x^{(i)}(s)\in C$ 的未移位次数。考虑由下式
定义的有序对
$(\widetilde x,\widetilde y)\in
\operatorname{Map}([0,1],C)^{d_1+\cdots+d_k+2-k}$：
\translatornote{原书把这个输出次数写成 $d_i+1$，其中 $i$ 在此并未固定，
且不可能给出 $k$ 元运算的总次数。由 $\deg\mathfrak m_k=2-k$ 以及
$x^{(i)}(s)$ 的未移位次数为 $d_i$，输出次数唯一确定为
$d_1+\cdots+d_k+2-k$；译文据此补正。}
\[
 \widetilde x(s)=\mathfrak m_k
 (x^{(1)}(s),\ldots,x^{(k)}(s)),
\]
并且当 $k\ne1$ 时，
\begin{equation}
 \widetilde y(s)=\sum_{j=1}^{k}
 (-1)^{d_1+\cdots+d_{j-1}+j}
 \mathfrak m_k(x^{(1)}(s),\ldots,y^{(j)}(s),\ldots,x^{(k)}(s));
\tag{4.2.12.1}\label{eq:4-2-12-1}
\end{equation}
当 $k=1$ 时，
\begin{equation}
 \widetilde y(s)=\frac{\partial x(s)}{\partial s}-\mathfrak m_1(y(s)).
\tag{4.2.12.2}\label{eq:4-2-12-2}
\end{equation}
于是定义
\begin{equation}
 \mathfrak m_k(\mathbf x^1,\ldots,\mathbf x^k)
 =(\widetilde x,\widetilde y),\qquad k\geq1.
\tag{4.2.12.3}\label{eq:4-2-12-3}
\end{equation}
令 $x_0=\mathfrak m_0(1)\in C$，其中 $\mathfrak m_0$ 是 $C$ 的
$A_\infty$ 结构的一部分。再以
\begin{equation}
 \mathfrak m_0(1)=(x_0,0)\in\operatorname{Map}([0,1],C)
\tag{4.2.12.4}\label{eq:4-2-12-4}
\end{equation}
定义 $\operatorname{Map}([0,1],C)$ 上的 $\mathfrak m_0$。
若 $\overline C$ 是 $A_\infty$ 代数，则以
\eqref{eq:4-2-12-1}--\eqref{eq:4-2-12-3} 的同样公式定义
$\overline{\operatorname{Map}}([0,1],\overline C)$ 上的
$\overline{\mathfrak m}_k$，并置 $\overline{\mathfrak m}_0=0$。

\begin{lemma}{4.2.13}\label{lem:4-2-13}
$(\operatorname{Map}([0,1],C),\mathfrak m)$ 是滤过 $A_\infty$ 代数，
而
$(\overline{\operatorname{Map}}([0,1],\overline C),
\overline{\mathfrak m})$ 是 $A_\infty$ 代数。若 $C$ 有严格幺元，则
$(\operatorname{Map}([0,1],C),\mathfrak m)$ 也有严格幺元；对
$\overline C$ 与
$(\overline{\operatorname{Map}}([0,1],\overline C),
\overline{\mathfrak m})$ 也有同样结论。
\end{lemma}

\begin{proof}
只证滤过 $A_\infty$ 代数的情形。需要计算
\[
 (\widehat x,\widehat y):=(\mathfrak m\circ\widehat d)
 (\mathbf x^{(1)},\ldots,\mathbf x^{(k)}).
\]
由 $(C,\mathfrak m)$ 的 $A_\infty$ 公式立即得到 $\widehat x=0$。
下面计算 $\widehat y=\widehat y(s)$。置
\[
 {}_{(j)}\mathbf x(s)=x^{(1)}(s)\otimes\cdots\otimes x^{(j-1)}(s),
 \qquad
 \mathbf x_{(j)}(s)=x^{(j+1)}(s)\otimes\cdots\otimes x^{(k)}(s),
\]
并记
\[
\begin{aligned}
 \Delta^{m-1}\mathbf x
  &=\sum_a\mathbf x_a^{(1)}\otimes\cdots\otimes\mathbf x_a^{(m)},\\
 \Delta^{m-1}{}_{(j)}\mathbf x
  &=\sum_a{}_{(j)}\mathbf x_a^{(1)}\otimes\cdots\otimes
    {}_{(j)}\mathbf x_a^{(m)},\\
 \Delta^{m-1}\mathbf x_{(j)}
  &=\sum_a\mathbf x_{(j)a}^{(1)}\otimes\cdots\otimes
    \mathbf x_{(j)a}^{(m)}.
\end{aligned}
\]
于是
{\small
\begin{equation}
\begin{aligned}
\widehat y(s)={}&
 \sum_a\sum_j(-1)^{\epsilon_1(a)+1}
 \mathfrak m\bigl({}_{(j)}\mathbf x_a^{(1)}(s),
   \mathfrak m({}_{(j)}\mathbf x_a^{(2)}(s)),
   {}_{(j)}\mathbf x_a^{(3)}(s),y^{(j)}(s),\mathbf x_{(j)}(s)\bigr)\\
&+\sum_a\sum_j(-1)^{\epsilon_2(a)+1}
 \mathfrak m\bigl({}_{(j)}\mathbf x(s),y^{(j)}(s),
   \mathbf x_{(j)a}^{(1)}(s),
   \mathfrak m(\mathbf x_{(j)a}^{(2)}(s)),
   \mathbf x_{(j)a}^{(3)}(s)\bigr)\\
&+\sum_a\sum_b\sum_j(-1)^{\epsilon_3(a)}
 \mathfrak m\bigl({}_{(j)}\mathbf x_a^{(1)}(s),
   \mathfrak m({}_{(j)}\mathbf x_a^{(2)}(s),y^{(j)}(s),
     \mathbf x_{(j)b}^{(1)}(s)),
   \mathbf x_{(j)b}^{(2)}(s)\bigr)\\
&-\sum_j\mathfrak m\left({}_{(j)}\mathbf x(s),
   \frac{\partial x^{(j)}(s)}{\partial s},\mathbf x_{(j)}(s)\right)
 +\frac{\partial}{\partial s}
   \mathfrak m_k(x^{(1)}(s),\ldots,x^{(k)}(s)),
\end{aligned}
\tag{4.2.14}\label{eq:4-2-14}
\end{equation}
}
其中
{\small
\[
\begin{aligned}
 \epsilon_1(a)
  &={}\deg'{}_{(j)}\mathbf x_a^{(1)}
     +\deg'{}_{(j)}\mathbf x_a^{(1)}
     +\deg'{}_{(j)}\mathbf x_a^{(2)}+1
     +\deg'{}_{(j)}\mathbf x_a^{(3)}\\
  &=\deg'{}_{(j)}\mathbf x
     +\deg'{}_{(j)}\mathbf x_a^{(1)}+1,\\
 \epsilon_2(a)
  &=\deg'{}_{(j)}\mathbf x+\deg y^{(j)}(s)+1
     +\deg'\mathbf x_{(j)a}^{(1)}+\deg'{}_{(j)}\mathbf x,\\
 \epsilon_3(a)
  &=\deg'{}_{(j)}\mathbf x_a^{(1)}
    +\deg'{}_{(j)}\mathbf x_a^{(2)}
    +\deg'\mathbf x_{(j)a}^{(1)}\\
  &=\deg'{}_{(j)}\mathbf x
    +\deg'{}_{(j)}\mathbf x_a^{(1)}.
\end{aligned}
\]
}
这里 $\deg'$ 是移位次数。由 $(C,\mathfrak m)$ 的 $A_\infty$ 公式，
\eqref{eq:4-2-14} 的前三行之和为零；其最后两行彼此抵消。因此
$A_\infty$ 公式成立。

若 $e$ 是 $C$ 的幺元，则 $(e,0)$ 是
$\operatorname{Map}([0,1],C)$ 的幺元。
\end{proof}

\begin{statement}{定义--命题}{4.2.15}\label{def-prop:4-2-15}
若 $R\supseteq\mathbb Q$，则
$(\operatorname{Map}([0,1],C),\mathfrak m)$ 是
$[0,1]\times C$ 的一个模型，而
$(\overline{\operatorname{Map}}([0,1],\overline C),
\overline{\mathfrak m})$ 是 $[0,1]\times\overline C$ 的一个模型。
\end{statement}

\begin{proof}
取 $s_0\in[0,1]$，定义 $A_\infty$ 同态
\[
 \mathrm{Eval}_{s=s_0}\colon\operatorname{Map}([0,1],C)\longrightarrow C
\]
如下：对 $\mathbf x=(x,y)$，置
$(\mathrm{Eval}_{s=s_0})_1(\mathbf x)=x(s_0)\in C$，并在 $k\ne1$ 时
置 $(\mathrm{Eval}_{s=s_0})_k=0$。

注意，\eqref{eq:4-2-11} 不是多项式，而是 $s$ 的形式幂级数。不过，仍可
代入 $x(s_0)$，因为正如注 4.2.10 所说，\eqref{eq:4-2-11} 的右端在
$\Lambda_{0,\mathrm{nov}}$ 的进拓扑中收敛。

容易验证 $\mathrm{Eval}_{s=s_0}$ 是滤过 $A_\infty$ 同态。再把 $C[1]$
的元素看作下列空间中的常值映射：
\[
 \operatorname{Map}([0,1],C[1]).
\]
以此定义
\[
 \mathrm{Incl}_1\colon C[1]\longrightarrow
 \operatorname{Map}([0,1],C[1]).
\]
若在 $k\ne1$ 时置 $\mathrm{Incl}_k=0$，便容易看出这给出滤过
$A_\infty$ 同态。

以同样方式定义 $\overline{\mathrm{Eval}}_{s=s_0}$ 与
$\overline{\mathrm{Incl}}$。条件 (4.2.3.1)、(4.2.3.2)、(4.2.3.4)、
(4.2.2.1)、(4.2.2.2)、(4.2.2.4) 都容易验证。因此，要证命题 4.2.15，
只需证明下一个引理。
\end{proof}

\begin{lemma}{4.2.16}\label{lem:4-2-16}
若 $R\supseteq\mathbb Q$，则 $\overline{\mathrm{Incl}}$ 在
$\overline{\mathfrak m}_1$ 上同调上诱导同构
\[
 H(\overline C,\overline{\mathfrak m}_1)
 \longrightarrow
 H(\overline{\operatorname{Map}}([0,1],\overline C),
   \overline{\mathfrak m}_1).
\]
\end{lemma}

\begin{proof}
$\overline{\mathrm{Incl}}_1$ 的像是
$(\overline{\operatorname{Map}}([0,1],\overline C),
\overline{\mathfrak m}_1)$ 的子复形；商复形
\[
 \overline{\operatorname{Map}}([0,1],\overline C)
 /\operatorname{Im}\overline{\mathrm{Incl}}_1
\]
同由形如 $\mathbf x(s)=(sx(s),y(s))$ 的元素组成的子复形同构。把这个
子复形记作
$(\overline{\operatorname{Map}}'([0,1],\overline C),
\overline{\mathfrak m}_1)$。

先考察
$(\overline{\operatorname{Poly}}'([0,1],\overline C),
\overline{\mathfrak m}_1)$；其他情形类似。有正合序列
\begin{center}
\resizebox{0.98\linewidth}{!}{$
0\longrightarrow(\overline C,\overline{\mathfrak m}_1)\otimes\mathbb Q[s]
\xrightarrow{\ i\ }
(\overline{\operatorname{Poly}}'([0,1],\overline C),
 \overline{\mathfrak m}_1)
\xrightarrow{\ \pi\ }
(\overline C,\overline{\mathfrak m}_1)\otimes\mathbb Q[s]
\longrightarrow0.
$}
\par\smallskip\textbf{图式 4.2.1}\label{diag:4-2-1}
\end{center}
这里 $\mathbb Q[s]$ 是多项式环，映射 $i$、$\pi$ 定义为
\[
 i(x(s))=(0,-x(s)),
 \qquad
 \pi(sx(s),y(s))=x(s).
\]
由 \eqref{eq:4-2-12-2}，图式 4.2.1 所诱导长正合序列的连接同态为
\begin{equation}
 x(s)\longmapsto-\frac{\partial}{\partial s}(sx(s)).
\tag{4.2.17}\label{eq:4-2-17}
\end{equation}
这个映射在 $\mathbb Q[s]$ 上是同构。这里正是需要假定
$\mathbb Q\subseteq R$ 的地方，因为 \eqref{eq:4-2-17} 的逆映射用到了
积分，即由
\[
 g\longmapsto\frac1s\int_0^s g(t)\,dt
\]
给出。这证明了 $\operatorname{Poly}$ 情形的引理；其余情形可用类似方式
证明。
\end{proof}

这样就在 $R\supseteq\mathbb Q$ 时构造出了 $[0,1]\times C$ 的一个模型。
模型 $\operatorname{Poly}([0,1],C)$ 适用于我们的绝大多数目的；但当
$R$ 有挠元时，它不能使用。因此下面在一般情形中另构造一个模型。本小节
余下部分不再假定 $R\supseteq\mathbb Q$。

置
\[
 C^{[0,1]}=C\oplus C[-1]\oplus C.
\]
定义 $\mathfrak I_0,\mathfrak I_1\colon C\to C^{[0,1]}$ 与
$\mathfrak J\colon C\to C^{[0,1]}$ 为
\[
 \mathfrak I_0(x)=(x,0,0),\qquad
 \mathfrak I_1(x)=(0,0,x),\qquad
 \mathfrak J(x)=(0,x,0).
\]
注意 $\mathfrak I_0,\mathfrak I_1$ 保持次数，而 $\mathfrak J$ 的次数为
$+1$。把 $\mathfrak I_0,\mathfrak I_1$ 延拓到
$BC[1]\to BC^{[0,1]}[1]$，仍用同样记号。我们将在 $C^{[0,1]}$ 上定义
一个 $A_\infty$ 结构，由此得到 $C$ 的另一个模型。在这个构造中，把
$C^{[0,1]}$ 上的滤过 $A_\infty$ 结构记作 $\mathfrak M_k$，把 $C$ 上的
滤过 $A_\infty$ 结构记作 $\mathfrak m_k$。定义
\begin{equation}
\left\{
\begin{aligned}
 \mathfrak M_1(\mathfrak I_0(x))
  &=\mathfrak I_0(\mathfrak m_1x)+(-1)^{\deg'x}\mathfrak J(x),\\
 \mathfrak M_1(\mathfrak I_1(x))
  &=\mathfrak I_1(\mathfrak m_1x)-(-1)^{\deg'x}\mathfrak J(x),\\
 \mathfrak M_1(\mathfrak J(x))&=\mathfrak J(\mathfrak m_1x),
\end{aligned}
\right.
\tag{4.2.18}\label{eq:4-2-18}
\end{equation}
以及
\begin{equation}
\left\{
\begin{aligned}
 (\mathrm{Eval}_{s=0})_1(x,y,z)&=x,
 & (\mathrm{Eval}_{s=1})_1(x,y,z)&=z,\\
 (\mathrm{Incl})_1(x)&=\mathfrak I_0(x)+\mathfrak I_1(x)=(x,0,x),
\end{aligned}
\right.
\tag{4.2.19}\label{eq:4-2-19}
\end{equation}
并置
\begin{equation}
 \mathfrak M_0(1)=(\mathrm{Incl})_1(\mathfrak m_0(1)).
\tag{4.2.20}\label{eq:4-2-20}
\end{equation}

\begin{remark}{4.2.21}\label{rem:4-2-21}
\begin{enumerate}[label=\textup{(\arabic*)},leftmargin=2.8em]
 \item $C^{[0,1]}$ 有如下几何来源。考虑 $C=C(L)$ 的情形。对
  $\sigma\in C(L)$，考察 $(-\epsilon,1+\epsilon)\times L$ 中形如
  \[
   (-\epsilon,1/2]\times\sigma,\qquad
   [1/2,1+\epsilon)\times\sigma,\qquad
   \{1/2\}\times\sigma
  \]
  的局部有限奇异链。这些链组成的链复形同 $C^{[0,1]}$ 同构；本注中置
  $\mathfrak m_0=0$。然而，若在这个几何模型中考察乘积结构（杯积），
  就会遇到困难，因为这些链彼此不横截；例如
  $\{1/2\}\times\sigma$ 不与 $\{1/2\}\times\sigma'$ 横截。为了定义
  乘积结构，必须以某种方式打破对称性，正如第 7 章将做的那样。下面的
  构造可以看作这一做法的代数类比。
 \item 上面写的是 $(-\epsilon,1+\epsilon)\times L$，而不是
  $L\times(-\epsilon,1+\epsilon)$；不过，\eqref{eq:4-2-18} 等式中
  所取的符号，是考察后一乘积时的符号。
\end{enumerate}
\end{remark}

现在对 $k\geq2$ 定义 $C^{[0,1]}$ 上的 $\mathfrak M_k$。只需在形如
$\mathfrak I_0(x)$、$\mathfrak I_1(x)$ 与 $\mathfrak J(x)$ 的元素上定义
这些算子。令
$\mathbf x\in B_kC$、$\mathbf z\in B_\ell C$、$y\in C$，其中
$k,\ell=0,1,2,\ldots$。定义
\begin{equation}
 \mathfrak M_{k+\ell+1}
 (\mathfrak I_0(\mathbf x),\mathfrak J(y),\mathfrak I_1(\mathbf z))
 =(-1)^{\deg'\mathbf z}\mathfrak J
 (\mathfrak m_{k+\ell+1}(\mathbf x,y,\mathbf z)),
\tag{4.2.22}\label{eq:4-2-22}
\end{equation}
并且对 $k,\ell\geq2$，定义
\begin{equation}
 \mathfrak M_k(\mathfrak I_0(\mathbf x))
 =\mathfrak I_0(\mathfrak m_k(\mathbf x)),
 \qquad
 \mathfrak M_\ell(\mathfrak I_1(\mathbf z))
 =\mathfrak I_1(\mathfrak m_\ell(\mathbf z)).
\tag{4.2.23}\label{eq:4-2-23}
\end{equation}
除 \eqref{eq:4-2-18}、\eqref{eq:4-2-20}、\eqref{eq:4-2-22}、
\eqref{eq:4-2-23} 给出的元素外，算子 $\mathfrak M_k$ 在其余元素上均取
零。也就是说，
\begin{align}
 \mathfrak M(\ldots,\mathfrak J(y_1),\ldots,
   \mathfrak J(y_2),\ldots)&=0,
 &\tag{4.2.24.1}\label{eq:4-2-24-1}\\
 \mathfrak M(\ldots,\mathfrak J(y_1),\ldots,
   \mathfrak I_0(x),\ldots)&=0,
 &\tag{4.2.24.2}\label{eq:4-2-24-2}\\
 \mathfrak M(\ldots,\mathfrak I_1(z),\ldots,
   \mathfrak J(y),\ldots)&=0,
 &\tag{4.2.24.3}\label{eq:4-2-24-3}\\
 \mathfrak M(\ldots,\mathfrak I_1(z),\ldots,
   \mathfrak I_0(x),\ldots)&=0,
 &\tag{4.2.24.4}\label{eq:4-2-24-4}\\
 \mathfrak M_{k+\ell}(\mathfrak I_0(\mathbf x),
   \mathfrak I_1(\mathbf z))&=0.
 &\tag{4.2.24.5}\label{eq:4-2-24-5}
\end{align}

\begin{lemma}{4.2.25}\label{lem:4-2-25}
$(C^{[0,1]},\mathfrak M_k)$ 是滤过 $A_\infty$ 代数。
\end{lemma}

\begin{proof}
检验 $A_\infty$ 关系 $\mathfrak M\circ\widehat d=0$；采用第 1 章末的
(Conv.3)。先计算
{\scriptsize
\[
\begin{aligned}
&(\mathfrak M\circ\widehat d)
 (\mathfrak I_0(\mathbf x),\mathfrak J(y),\mathfrak I_1(\mathbf z))\\
={}&\sum_{c_1,c_2}(-1)^{
 \deg'\mathbf x_{c_1}^{(2;1)}+\deg'\mathbf z_{c_2}^{(2;1)}}
 \mathfrak M\bigl(\mathfrak I_0(\mathbf x_{c_1}^{(2;1)}),
  \mathfrak J(\mathfrak m(\mathbf x_{c_1}^{(2;2)},y,
    \mathbf z_{c_2}^{(2;1)})),
  \mathfrak I_1(\mathbf z_{c_2}^{(2;2)})\bigr)\\
&+\sum_c(-1)^{\deg'\mathbf x_c^{(3;1)}}
 \mathfrak M\bigl(\mathfrak I_0(\mathbf x_c^{(3;1)}),
  \mathfrak M(\mathfrak I_0(\mathbf x_c^{(3;2)})),
  \mathfrak I_0(\mathbf x_c^{(3;3)}),\mathfrak J(y),
  \mathfrak I_1(\mathbf z)\bigr)\\
&+\sum_c(-1)^{\deg'\mathbf z_c^{(3;1)}+
 \deg'\mathbf x+\deg'y+1}
 \mathfrak M\bigl(\mathfrak I_0(\mathbf x),\mathfrak J(y),
  \mathfrak I_1(\mathbf z_c^{(3;1)}),
  \mathfrak M(\mathfrak I_1(\mathbf z_c^{(3;2)})),
  \mathfrak I_1(\mathbf z_c^{(3;3)})\bigr).
\end{aligned}
\]
}
最后一行 $(-1)$ 的指数中，$+1$ 代表 $\deg\mathfrak J=+1$。因此得到
{\scriptsize
\[
\begin{aligned}
&(\mathfrak M\circ\widehat d)
 (\mathfrak I_0(\mathbf x),\mathfrak J(y),\mathfrak I_1(\mathbf z))\\
={}&\sum_{c_1,c_2}(-1)^{
 \deg'\mathbf x_{c_1}^{(2;1)}+\deg'\mathbf z}
 \mathfrak J\circ\mathfrak m
 (\mathbf x_{c_1}^{(2;1)},
  \mathfrak m(\mathbf x_{c_1}^{(2;2)},y,\mathbf z_{c_2}^{(2;1)}),
  \mathbf z_{c_2}^{(2;2)})\\
&+\sum_c(-1)^{\deg'\mathbf x_c^{(3;1)}+\deg'\mathbf z}
 \mathfrak J\circ\mathfrak m
 (\mathbf x_c^{(3;1)},\mathfrak m(\mathbf x_c^{(3;2)}),
  \mathbf x_c^{(3;3)},y,\mathbf z)\\
&+\sum_c(-1)^{\deg'\mathbf z_c^{(3;1)}+
 \deg'\mathbf x+\deg'y-1+\deg'\mathbf z+1}
 \mathfrak J\circ\mathfrak m
 (\mathbf x,y,\mathbf z_c^{(3;1)},
  \mathfrak m(\mathbf z_c^{(3;2)}),\mathbf z_c^{(3;3)}).
\end{aligned}
\]
}
最后一行中，
\[
 \deg'\mathbf z+1
 =\deg(\mathbf z_c^{(3;1)},\mathfrak m(\mathbf z_c^{(3;2)}),
       \mathbf z_c^{(3;3)}).
\]
故上式由 $C$ 上的 $A_\infty$ 关系等于零。

\begin{remark}{4.2.26}\label{rem:4-2-26}
上面第一个等式第三行中的
$\mathfrak M(\mathfrak I_0(\mathbf x_c^{(3;2)}))$ 包含
$\mathbf x_c^{(3;2)}\in B_1C$ 或 $B_0(C)$ 的情形。在这些情形中，
\[
 \mathfrak M(\mathfrak I_0(\mathbf x_c^{(3;2)}))
 \ne\mathfrak I_0(\mathfrak m(\mathbf x_c^{(3;2)})).
\]
不过，由 (4.2.24)，含有
\[
 \mathfrak M(\mathfrak I_0(\mathbf x_c^{(3;2)}))
 -\mathfrak I_0(\mathfrak m(\mathbf x_c^{(3;2)}))
\]
的项消失。
\end{remark}

再计算 $(\mathfrak M\circ\widehat d)(\mathfrak I_0(\mathbf x))$。写
$\mathbf x=\mathbf x_-\otimes x_{\mathrm{last}}$，其中
$x_{\mathrm{last}}\in B_1C[1]$，则
{\small
\begin{equation}
\begin{aligned}
&(\mathfrak M\circ\widehat d)(\mathfrak I_0(\mathbf x))\\
={}&\sum_c(-1)^{\deg'\mathbf x_c^{(3;1)}}
 \mathfrak M\bigl(\mathfrak I_0(\mathbf x_c^{(3;1)}),
  \mathfrak I_0(\mathfrak m(\mathbf x_c^{(3;2)})),
  \mathfrak I_0(\mathbf x_c^{(3;3)})\bigr)\\
&+(-1)^{\deg'\mathbf x_-+\deg'x_{\mathrm{last}}}
 \mathfrak M(\mathfrak I_0(\mathbf x_-),\mathfrak J(x_{\mathrm{last}}))\\
={}&\sum_c(-1)^{\deg'\mathbf x_c^{(3;1)}}
 \mathfrak I_0\bigl(\mathfrak m(\mathbf x_c^{(3;1)},
  \mathfrak m(\mathbf x_c^{(3;2)}),\mathbf x_c^{(3;3)})\bigr)\\
&+(-1)^{\deg'\mathbf x_-+\deg'x_{\mathrm{last}}}
 \mathfrak J(\mathfrak m(\mathbf x_-,x_{\mathrm{last}}))
 +(-1)^{\deg'\mathfrak m(\mathbf x)}
 \mathfrak J(\mathfrak m(\mathbf x)).
\end{aligned}
\tag{4.2.27}\label{eq:4-2-27}
\end{equation}
}
最后一行来自第一行中 $\mathbf x_c^{(3;2)}=\mathbf x$ 的情形。末两项
彼此抵消，故 \eqref{eq:4-2-27} 为零。

再计算 $(\mathfrak M\circ\widehat d)(\mathfrak I_1(\mathbf z))$。写
$\mathbf z=z_{\mathrm{first}}\otimes\mathbf z_-$，其中
$z_{\mathrm{first}}\in B_1C[1]$，则
{\small
\begin{equation}
\begin{aligned}
&(\mathfrak M\circ\widehat d)(\mathfrak I_1(\mathbf z))\\
={}&\sum_c(-1)^{\deg'\mathbf z_c^{(3;1)}}
 \mathfrak M\bigl(\mathfrak I_1(\mathbf z_c^{(3;1)}),
  \mathfrak I_1(\mathfrak m(\mathbf z_c^{(3;2)})),
  \mathfrak I_1(\mathbf z_c^{(3;3)})\bigr)\\
&-(-1)^{\deg'z_{\mathrm{first}}}
 \mathfrak M(\mathfrak J(z_{\mathrm{first}}),\mathfrak I_1(\mathbf z_-))\\
={}&\sum_c(-1)^{\deg'\mathbf z_c^{(3;1)}}
 \mathfrak I_1\bigl(\mathfrak m(\mathbf z_c^{(3;1)},
  \mathfrak m(\mathbf z_c^{(3;2)}),\mathbf z_c^{(3;3)})\bigr)\\
&-(-1)^{\deg'z_{\mathrm{first}}+\deg'\mathbf z_-}
 \mathfrak J(\mathfrak m(z_{\mathrm{first}},\mathbf z_-))
 -(-1)^{\deg'\mathfrak m(\mathbf z)}
 \mathfrak J(\mathfrak m(\mathbf z)).
\end{aligned}
\tag{4.2.28}\label{eq:4-2-28}
\end{equation}
}
最后一行来自第一行中 $\mathbf z_c^{(3;2)}=\mathbf z$ 的情形。末两项
彼此抵消，故 \eqref{eq:4-2-28} 为零。

最后，对 $\mathbf x\in B_kC[1]$、$\mathbf z\in B_\ell C[1]$，其中
$k,\ell>0$，计算
$(\mathfrak M\circ\widehat d)
(\mathfrak I_0(\mathbf x),\mathfrak I_1(\mathbf z))$。有下式。
\translatornote{原书在 \eqref{eq:4-2-28} 后把
$\mathbf z_c^{(3;2)}=\mathbf z$ 误写成
$\mathbf x_c^{(3;2)}=\mathbf x$；又在 \eqref{eq:4-2-29} 第二项的
$\mathfrak J(z_{\mathrm{first}})$ 与 $\mathfrak I_1(\mathbf z_-)$ 之间
漏写逗号。两处分别由正在计算的输入及已定义的三元输入型别唯一确定，译文
据此改正。}
{\small
\begin{equation}
\begin{aligned}
&(\mathfrak M\circ\widehat d)
 (\mathfrak I_0(\mathbf x),\mathfrak I_1(\mathbf z))\\
={}&(-1)^{\deg'\mathbf x_-+\deg'x_{\mathrm{last}}}
 \mathfrak M(\mathfrak I_0(\mathbf x_-),\mathfrak J(x_{\mathrm{last}}),
             \mathfrak I_1(\mathbf z))\\
&-(-1)^{\deg'\mathbf x+\deg'z_{\mathrm{first}}}
 \mathfrak M(\mathfrak I_0(\mathbf x),\mathfrak J(z_{\mathrm{first}}),
             \mathfrak I_1(\mathbf z_-))\\
={}&(-1)^{\deg'\mathbf x_-+\deg'x_{\mathrm{last}}+\deg'\mathbf z}
 \mathfrak J(\mathfrak m(\mathbf x_-,x_{\mathrm{last}},\mathbf z))\\
&-(-1)^{\deg'\mathbf x+\deg'z_{\mathrm{first}}+\deg'\mathbf z_-}
 \mathfrak J(\mathfrak m(\mathbf x,z_{\mathrm{first}},\mathbf z_-))=0.
\end{aligned}
\tag{4.2.29}\label{eq:4-2-29}
\end{equation}
}
\end{proof}

\begin{statement}{定义--引理}{4.2.30}\label{def-lem:4-2-30}
$(C^{[0,1]},\mathfrak M_k)$\liftIndex{I144} 是 $[0,1]\times C$ 的一个模型。若 $C$ 有
幺元，则 $(C^{[0,1]},\mathfrak M_k)$ 也有幺元。
\end{statement}

\begin{proof}
$\mathrm{Eval}_{s=i}$ 与 $\mathrm{Incl}$ 已经定义；容易看出它们诱导上链
同伦等价。若 $e$ 是 $C$ 的幺元，则 $(e,0,e)$ 是
$(C^{[0,1]},\mathfrak M_k)$ 的幺元。
\end{proof}

以同样方式可以定义
$(\overline C^{[0,1]},\overline{\mathfrak M}_k)$；它是
$[0,1]\times\overline C$ 的一个模型。由此已经证明：

\begin{theorem}{4.2.31}\label{thm:4-2-31}
$[0,1]\times\overline C$ 的模型 $\overline{\mathfrak C}$ 以及
$[0,1]\times C$ 的模型 $\mathfrak C$ 总是存在。若 $C$ 有隙，则可取
$\mathfrak C$ 也有隙。若 $\overline C$（相应地，$C$）有幺元，则可取
$\overline{\mathfrak C}$（相应地，$\mathfrak C$）为含幺模型。
\end{theorem}

\begin{remark}{4.2.32}\label{rem:4-2-32}
上面通过直接定义运算来定义 $C^{[0,1]}$。也可以利用障碍理论归纳证明这些
运算存在，从而证明定理 4.2.31。
\end{remark}

\begin{remark}{4.2.33}\label{rem:4-2-33}
$A_\infty$ 代数 $\overline{\operatorname{Poly}}([0,1],\overline C)$ 在
多方面优于 $\overline C^{[0,1]}$。例如，若 $\overline C$ 是（分次）交换
D.G.A.，则 $\overline{\operatorname{Poly}}([0,1],\overline C)$ 也是分次
交换 D.G.A.；另一方面，$\overline C^{[0,1]}$ 总是非交换的。作者不知道
如何在挠系数上构造保持交换性的 $[0,1]\times\overline C$ 模型。

$\overline C^{[0,1]}$ 同沙利文
在 D.G.A. 语境中使用的模型 \textbf{[Sul78]} 有些相似。不过，沙利文所用
模型是交换的，而且只适用于极小 D.G.A. 的情形。
\end{remark}

下面这个关键结果用于建立（滤过）$A_\infty$ 同态之间同伦的基本性质。

\begin{theorem}{4.2.34}\label{thm:4-2-34}
令 $C_1,C_2$ 为有隙滤过 $A_\infty$ 代数，令
$\mathfrak C_1,\mathfrak C_2$ 分别是 $[0,1]\times C_1$、
$[0,1]\times C_2$ 的任意有隙模型。再令
$\mathfrak f\colon C_1\to C_2$ 为有隙滤过 $A_\infty$ 同态。那么存在
有隙滤过 $A_\infty$ 同态
$\mathfrak F\colon\mathfrak C_1\to\mathfrak C_2$，使
\[
 \mathrm{Eval}_{s=s_0}\circ\mathfrak F
 =\mathfrak f\circ\mathrm{Eval}_{s=s_0},
 \qquad s_0=0,1,
\]
且
\[
 \mathrm{Incl}\circ\mathfrak f=\mathfrak F\circ\mathrm{Incl}.
\]
无滤过版本也成立。若 $\mathfrak f$ 含幺，则 $\mathfrak F$ 也含幺；若
$\mathfrak f$ 严格，则 $\mathfrak F$ 也严格。
\end{theorem}

\begin{center}
\resizebox{0.88\linewidth}{!}{$
\begin{array}{ccccc}
C_1&\xrightarrow{\ \mathrm{Incl}\ }&\mathfrak C_1
&\xrightarrow{\ \mathrm{Eval}_{s=0}\oplus\mathrm{Eval}_{s=1}\ }&C_1\oplus C_1\\
{\scriptstyle\mathfrak f}\big\downarrow&&
{\scriptstyle\mathfrak F}\big\downarrow&&
\big\downarrow{\scriptstyle\mathfrak f\oplus\mathfrak f}\\
C_2&\xrightarrow{\ \mathrm{Incl}\ }&\mathfrak C_2
&\xrightarrow{\ \mathrm{Eval}_{s=0}\oplus\mathrm{Eval}_{s=1}\ }&C_2\oplus C_2.
\end{array}
$}
\par\smallskip\textbf{图式 4.2.2}\label{diag:4-2-2}
\end{center}

若上述两个关系成立，就称 $\mathfrak F$ 在
$\mathfrak f\colon C_1\to C_2$ 之上与 $\mathrm{Incl}$、
$\mathrm{Eval}$ 相容。定理 4.2.34 将在第 4.4 节证明。
